% !TEX encoding = Windows Latin 1


%
% Legt die Grundstruktur bzw. die Art des Dokuments sowie die Standardschriftgr"a"ae fest.
%
\documentclass[11pt]{article}

%
% Legt das Seitenlayout fest.
%
\usepackage[a4paper, margin=1.0in]{geometry}

%
% Legt die verwendete Schrift fest.
%
\usepackage{times}

%
% Legt die Worttrennungsregeln fest.
%
\usepackage[ngerman]{babel}

%
% Legt die Eingabekodierung fest.
%
\usepackage[latin1]{inputenc}
\usepackage[T1]{fontenc}

%
% Erm"aglicht die Darstellung von Farben sowie die Bezeichnung von Farben durch Namen.
%
\usepackage[usenames, dvipsnames]{color}

%
% Erm"aglicht das Einf"agen von Grafiken.
%
\usepackage{graphicx}

%
% Erleichtert die Darstellung von URLs.
%
\usepackage{url}

%
% Erlaubt die Bearbeitung von Bilder-, Diagramm- und Tabellenunterschriften.
%
\usepackage[hang, normalsize, bf]{caption}

%
% Legt fest, dass die AMS-Symbole benutzt werden.
%
\usepackage{amsmath}
\usepackage{amssymb}

%
% Erm"aglicht die Verwendung von vielen mathematischen Symbolen.
%
\usepackage{mathabx}

%
% Erm"aglicht die Verwendung der \today-Anweisung.
%
\usepackage[nodayofweek]{datetime}

% Erleichtert die Formattierung von Quelltexten.
\usepackage{listings}

%
% Erm"aglicht die einfachere Bearbeitung von Tabellen.
%
\usepackage{tabularx}

%
% Erm"aglicht die Darstellung von Tabellen auf mehreren Seiten (s. Dokumentation).
%
\usepackage{ltxtable}

%
% Erm"aglicht, Appendices in das Inhaltverzeichnis zu "abernehmen.
%
\usepackage[toc]{appendix}

%
% Erm"aglicht das Zusammenf"agen mehrerer Zeilen einer Tabelle.
%
\usepackage{multirow}

%
% Legt fest, welche Sprache in der Bibliographie verwendet werden soll.
%
%\usepackage[fixlanguage]{babelbib}

%
% Legt PDF-Format: PDF-A-1a fest
%
\usepackage[a-1a]{pdfx}

%%%%%%%%%%%%%%%%%%%%%%%%%%%%%%%%%%%%%%%%%%%%%%%%%%%%%%%%%%%%%%%%%%%%%%%%%%%%%%%%
%                                                                              %
% Einstellungen                                                                %
%                                                                              %
%%%%%%%%%%%%%%%%%%%%%%%%%%%%%%%%%%%%%%%%%%%%%%%%%%%%%%%%%%%%%%%%%%%%%%%%%%%%%%%%

%%%%%%%%%%%%%%%%%%%%%%%%%%%%%%%%%%%%%%%%%%%%%%%%%%%%%%%%%%%%%%%%%%%%%%%%%%%%%%%%
% einstellungen: url
%%%%%%%%%%%%%%%%%%%%%%%%%%%%%%%%%%%%%%%%%%%%%%%%%%%%%%%%%%%%%%%%%%%%%%%%%%%%%%%%

\urlstyle{tt}

%%%%%%%%%%%%%%%%%%%%%%%%%%%%%%%%%%%%%%%%%%%%%%%%%%%%%%%%%%%%%%%%%%%%%%%%%%%%%%%%
% einstellungen: bibliographie
%%%%%%%%%%%%%%%%%%%%%%%%%%%%%%%%%%%%%%%%%%%%%%%%%%%%%%%%%%%%%%%%%%%%%%%%%%%%%%%%

%\selectbiblanguage{german}
%\bibliographystyle{babplain}
\bibliographystyle{alpha}

%%%%%%%%%%%%%%%%%%%%%%%%%%%%%%%%%%%%%%%%%%%%%%%%%%%%%%%%%%%%%%%%%%%%%%%%%%%%%%%%
%                                                                              %
% ANWEISUNGEN                                                                  %
%                                                                              %
%%%%%%%%%%%%%%%%%%%%%%%%%%%%%%%%%%%%%%%%%%%%%%%%%%%%%%%%%%%%%%%%%%%%%%%%%%%%%%%%

\setlength{\parskip}{1.5ex plus 0.5ex minus 0.5ex}
\setlength{\parindent}{0em}

%%%%%%%%%%%%%%%%%%%%%%%%%%%%%%%%%%%%%%%%%%%%%%%%%%%%%%%%%%%%%%%%%%%%%%%%%%%%%%%%
% anweisungen: r"amische zahlen
%%%%%%%%%%%%%%%%%%%%%%%%%%%%%%%%%%%%%%%%%%%%%%%%%%%%%%%%%%%%%%%%%%%%%%%%%%%%%%%%

\makeatletter
\newcommand{\rmnum}[1]{\romannumeral #1}
\newcommand{\Rmnum}[1]{\expandafter\@slowromancap\romannumeral #1@}
\makeatother

%%%%%%%%%%%%%%%%%%%%%%%%%%%%%%%%%%%%%%%%%%%%%%%%%%%%%%%%%%%%%%%%%%%%%%%%%%%%%%%%
%                                                                              %
% DOKUMENT                                                                     %
%                                                                              %
%%%%%%%%%%%%%%%%%%%%%%%%%%%%%%%%%%%%%%%%%%%%%%%%%%%%%%%%%%%%%%%%%%%%%%%%%%%%%%%%


\begin{document}

\inputencoding{latin1}


%%%%%%%%%%%%%%%%%%%%%%%%%%%%%%%%%%%%%%%%%%%%%%%%%%%%%%%%%%%%%%%%%%%%%%%%%%%%%%%%
% dokument: titelblatt
%%%%%%%%%%%%%%%%%%%%%%%%%%%%%%%%%%%%%%%%%%%%%%%%%%%%%%%%%%%%%%%%%%%%%%%%%%%%%%%%

\begin{titlepage}

  %
  % university logo
  %

  \begin{flushright}
    \includegraphics[scale=0.75, keepaspectratio=true]{data/fig/logo.pdf}
  \end{flushright}

  %
  % university, department & lecture info
  %

  \begin{center}
    \ \\
    \textsc{\Large
      Universit"at Potsdam, \\Institut f"ur Informatik \\
      \ \\
      \ \\
      Software Engineering, WiSe 2017-18 \\
      \ \\
      \ \\
      Projektarbeit -- Teil I \\
      \ \\
      \ \\
    }
    {\huge \bfseries
      SharedBox Ultimate \\
    }
    \ \\
    \ \\
    \ \\
    \ \\
    \ \\
    \ \\
    \ \\

    %
    % author & supervisors
    %

    \large
    \begin{tabularx}{\textwidth}{ l X X }
      \emph{Autoren}:	&                   &	\emph{Betreuer}:        \\
                  	&                   &                           \\
      Vorname Name	& (Mat-Nr.\ 123456) &	Prof. Dr.\ Christian Hammer \\
      Vorname Name	& (Mat-Nr.\ 123456) &	M. Sc.  Sascha Gross \\
      Vorname Name	& (Mat-Nr.\ 123456) &	G�nther Wullaert  \\
      Vorname Name	& (Mat-Nr.\ 123456) &	Juliane Scherlitzki    \\
      Vorname Name	& (Mat-Nr.\ 123456) &	 \\
    \end{tabularx}

    \vfill
    {\Large \today}
  \end{center}
\end{titlepage}

%%%%%%%%%%%%%%%%%%%%%%%%%%%%%%%%%%%%%%%%%%%%%%%%%%%%%%%%%%%%%%%%%%%%%%%%%%%%%%%%
% dokument: inhaltsverzeichnis
%%%%%%%%%%%%%%%%%%%%%%%%%%%%%%%%%%%%%%%%%%%%%%%%%%%%%%%%%%%%%%%%%%%%%%%%%%%%%%%%

\tableofcontents{}
\newpage{}

%%%%%%%%%%%%%%%%%%%%%%%%%%%%%%%%%%%%%%%%%%%%%%%%%%%%%%%%%%%%%%%%%%%%%%%%%%%%%%%%
% dokument: kapitel
%%%%%%%%%%%%%%%%%%%%%%%%%%%%%%%%%%%%%%%%%%%%%%%%%%%%%%%%%%%%%%%%%%%%%%%%%%%%%%%%

\section{Daten}
In diesem Abschnitt sollen die wichtigsten Daten zu dem Projekt dargestellt
werden. Dazu z"ahlen Beginn und Ende des Projektes, beteiligte Personen
innerhalb des Projekt-Teams und beteiligte Personen au"serhalb des Projekt-Teams.

\subsection{Teilnehmer}
Folgende Tabelle kann als Vorlage verwendet werden:

\begin{tabular}{|c|c|c|c|c|}
\hline
\textbf{Name} &\textbf{Vorname} & \textbf{Position} &\textbf{intern/extern} &
\textbf{Kontakt (E-Mail)}\\\hline
Jenkins & Joe     & Team-Leader            & intern &\\\hline
Kloth   & Michael & Modellierer            & intern &\\\hline
Setsche & Gregor  & Programmierer          & intern &\\\hline
Gruber  & Stephan & Projekt-Manager        & intern &\\\hline
M"uller & Anton   & Auftraggeber           & extern &\\\hline
Mendel  & Juliane & Leitende Endanwenderin & extern &\\\hline
\end{tabular}

Externe beteiligte Personen k"onnen wichtig sein, wenn konkrete Absprachen
bez"uglich der Benutzung (User Stories) gekl"art werden m"ussen.

Wichtig ist, zu beachten, dass pers"onliche Daten vertraulich behandelt werden.

\section{Scope Management}
Das Scopemanagement besch"aftigt sich mit der Bestimmung welche
Arbeit zum und welche Arbeit nicht zum Projekt geh"ort. So l"asst sich der
Prozess in folgende Hauptprozesse \cite{SCHW09} unterteilen:
\begin{itemize}
  \item Bestimmung aller Anforderungen,
  \item Definition der Aufgabendefinition,
  \item Work Breakdown Structure erstellen,
  \item Pr"ufen und Kontrolle der Aufgaben.
\end{itemize}


\subsection{Aufgabendefinition}
Zur Bestimmung der Anforderungen ist eine Analyse des Anforderungsdokuments
n"otig. Aus dieser sollte hervorgehen welche Aufgaben gemacht werden m"ussen, bis wann diese gemacht werden m"ussen und zu welchem Preis. Der Preis spielt
hier in diesem speziellen Projekt keine Rolle. Die o.g. Aufgaben sollten so
fein wie m"oglich aufgeteilt werden, um sp"ater eine gute Aufgabenverteilung
zu erreichen. Im letzten Schritt werden die erstellten Anforderungen "uberpr"uft
und die laufende Arbeit bez"uglich dieser kontrolliert.

\section{Qualit"atsmanagement}
In diesem Abschnitt wird beschrieben, wie die Qualit"at der Projektergebnisse
sichergestellt wurde. Wichtige Aspekte sind dabei der Grad der Konformit"at zu
den Anforderungen und die Eignung zur Nutzung. Dabei ist zu beschreiben, wie
sichergestellt wurde, dass die Projekt-Prozesse und die (Teil-) Ergebnisse, die
durch die Anforderung beschriebene Spezifikation erf"ullen. Beziehungsweise wie
sichergestellt wurde, dass das Produkt so wie vorgesehen genutzt werden kann.

Im Qualit"atsmanagement gibt es unter Anderem folgende wichtige Prozesse, die
vom Manager ausgef"uhrt werden m"ussen:
\begin{itemize}
  \item Qualit"at Planung
  \item Qualit"at Sicherung
  \item Qualit"at Kontrolle
\end{itemize}

\subsection{Qualit"at Planung}
Hier sollen die f"ur das Projekt relevanten Qualit"ats-Standards identifiziert
werden. Dazu z"ahlt zum Beispiel die "Ubersichtlichkeit und Intuitivit"at von
Diagrammen.

\subsection{Qualit"at Sicherung}
Hier soll beschrieben werden, wie die Qualit"at der Ergebnisse gesichert wurde.
Beispielsweise k"onnten die Teilergebnisse auf die Erf"ullung der
Qualit"ats-Eigenschaften in konkreten Abst"anden "uberpr"uft worden sein.

Wichtige Aspekte der Qualit"at sind die Funktionalit"at (Wurden die vorgesehenen
Funktionen integriert?), die Features (Wurden W"unsche von Endanwendern, die die
Benutzbarkeit betreffen erf"ullt?), Outputs (Sind Ausgaben, die das System macht
f"ur einen Benutzer (intuitiv) verst"andlich?) , Performanz
(Wie gut/schnell/effizient f"uhrt das System die Prozesse aus?),
Zuverl"assigkeit (Arbeitet das Produkt bei gew"ohnlicher Nutzung einwandfrei)
und Wartbarkeit.

Wichtig ist nat"urlich auch, ob aufgrund nicht erf"ullter
Qualit"ats-Eigenschaften das Projekt angepasst oder Teile erneut erarbeitet
werden mussten und welche Ursachen zu der ungen"ugenden Qualit"at f"uhrten.

\subsection{Qualit"at Kontrolle}
Wurden die Sicherungsma"snahmen auch ausgef"uhrt und in welcher Weise wurde dies
dokumentiert?

\section{Personal Management}
Das Personalmanagement umfasst die Prozesse, die erforderlich sind, um einen
m"oglichst effektiven Einsatz der am Projekt beteiligten Mitarbeiter zu
organisieren.

Hauptprozesse \cite{SCHW09} sind:
\begin{itemize}
  \item Erstellung eines Planes zur Mitarbeiterentwicklung,
  \item Zusammenstellung des Projektteams,
  \item Entwicklung des Projektteams,
  \item Verwaltung des Projektteams.
\end{itemize}

\subsection{Aufgabenverteilung}
In diesem Abschnitt besch"aftigt man sich mit der Verwaltung der Mitarbeiter,
die zur Verf"ugung standen und wie und aus welchem Grund sie in der
konkreten Art und Weise eingesetzt wurden.

Hierbei ist zu beschreiben, welche Projekt-Mitglieder wof"ur eingesetzt wurden
und wie hoch die Arbeitslast geplant war. Ein wichtiger Gesichtspunkt sind die
Fertigkeiten und die Erfahrung, die die Projekt-Mitglieder haben und wie diese
genutzt wurden, um das Projekt fertigzustellen.
Mussten die Mitglieder eingearbeitet werden? Inwieweit wurde dies vom Management
unterst"utzt?

Es gibt Personen, die aufgrund ihres Auftretens auf andere Mitarbeiter Vertrauen
und F"uhrungsqualit"at/Autorit"at ausstrahlen. Wurden solche Personen genutzt,
um den Gruppenzusammenhalt zu f"ordern?

Wurden Konzepte zur Motivation angewandt, um den Arbeitswillen zu unterst"utzen?
Wie wurde verhindert, dass einzelne Projekt-Mitglieder "uberlastet werden?

Wurden Teilgruppen gebildet, um das Projekt in kleinere Teile aufl"osen zu
k"onnen? Wie wurde sichergestellt, dass die Zusammenarbeit zwischen den
Teilgruppen funktioniert?

Mussten w"ahrend der Laufzeit des Projektes weitere Mitarbeiter in das
Projekt-Team integriert werden. Wie wurde verhindert, dass die Synergien
innerhalb der Gruppe gest"ort wurden?

\section{Time Management}
Das Time Management besch"aftigt sich mit dem zeitlichen Ablauf eines Projektes.
Es beschreibt die Prozesse, welche notwendig sind, um das Projekt p"unktlich zu
liefern. Dazu geh"oren folgende Hauptprozesse \cite{SCHW09}:
\begin{itemize}
  \item Definition der Aktivit"aten,
  \item Bestimmung der Abh"angigkeiten zwischen den Aktivit"aten,
  \item Absch"atzung der Ressourcen f"ur jede Aktivit"at,
  \item Absch"atzung der Dauer der Aktivit"aten,
  \item Entwurf eines Zeitplans,
  \item Kontrolle des Zeitplans.
\end{itemize}
Die Definition von Aktivit"aten bezieht sich auf Erstellung von Teilen (nicht
nur Teilsystemen) der Gesamtaufgabe. Als n"achsten Schritt m"ussen zeitliche
Abh"angigkeiten
zwischen den Aktivit"aten "uberpr"uft und im Nachhinein bei der Erstellung des
Zeitplans ber"ucksichtigt werden. Als n"achsten Schritt macht man eine
Aufwandssch"atzung zur Bestimmung der ben"otigten Ressourcen (z.B. Arbeiter,
Equipment) und der Aktivit"atsdauer. Stellen Sie einen Bezug zu den Ressourcen
her. Auf der Basis der vorher beschriebenen Voraussetzungen erstellt man einen
Zeitplan, welcher alle Anforderung erf"ullt. Hierf"ur werden oft zeitliche
Diagramme in Form von Tabellen oder Gantt-Charts erstellt. Als letztes bleibt
nur noch die laufende "Uberpr"ufung der Einhaltung von Meilensteinen und ggf.
eine Anpassung des Zeitplans.

\subsection{Aufwandssch"atzung}

\subsection{Projektplan}

\subsection{Soll/Ist Vergleich}

\section{Besonderheiten}

\subsection{Beziehungen zwischen den Aufgabenteilen}
In diesem Kapitel werden die Besonderheiten des Projekts festgehalten. Der
Projektleiter soll dokumentieren, welche Auff"alligkeiten bemerkbar geworden
sind. Dazu z"ahlen die Zusammenh"ange in dem Projekt zwischen den einzelnen
Aufgaben. Wenn eine Aufgabe mit einer anderen sich sehr gut kombinieren l"asst,
oder sogar bei der Anfertigung Synergie Effekte entstehen (die Qualit"at der
nachfolgenden Aufgabe verbessert sich erheblich, oder der Zeitaufwand wird
geringer), dann geh"oren diese Beobachtungen in diesen Abschnitt.

Desweiteren ist es ratsam die Auff"alligkeiten in diesem Kapitel zu beschreiben
. Es ist wichtig zu identifizieren, wie sich das Verhalten der Auftraggeber auf
den Projektverlauf widergespiegelt hat. Wenn diese zum Beispiel nur schwammige
Antworten gegeben haben, welche die Arbeit nur verz"ogert haben, soll dieses
festgehalten werden.

\subsection{Externe Verz"ogerungen}
Externe Faktoren geh"oren ebenso in diesen Abschnitt. Unter externen Faktoren
kann die Einflussnahme von h"oherer Gewalt (Wetterbedingungen, Stromunterbrechung
usw.) mit einbezogen werden.  Eine Intervention seitens des Auftraggeber oder
dritter, welche erheblichen Einfluss auf die Projektplanung und den Verlauf haben,
kann hier beschrieben werden.

Die Besonderheiten der politischen und ethischen Seite des Projekts werden
aufgenommen. Es k"onnen Projekte im milit"arischen oder medizinischen Bereich
stattfinden, welche besondere Auflagen seitens des Gesetzgebers mit sich bringen
und wodurch au"serordentliche Rahmenbedingungen entstehen.

\section{Organisation}

\subsection{Abrechnung des Time- und Ressourcenmanagements}
In diesem Abschnitt werden die organisatorischen Einzelheiten des Projekts
festgehalten. In erster Linie ist das Verhalten der Teilnehmer zu beschreiben.
Die Fragen: Wie p"unktlich sind die Mitarbeiter? Wie oft gab es Verz"ogerungen?
Wurden die Aufgaben in angemessener Qualit"at erledigt? sind hier zu beantworten.
In diesem Zusammenhang kann kurz auf die Struktur der Aufgabenpakete eingegangen
werden, die Aufteilung der einzelnen Aufgaben und das Ma"s der Produktivit"at
definiert werden. Der angefertigte Stundenzettel der Mitarbeitern dient hier
als Grundlage.

\subsubsection{Stundenzettel}
Angaben zu den Arbeitszetteln von den Mitarbeitern sowie Aufnahme der
Arbeitsaufteilung. Beschreibung der ausgew"ahlten Methode zur Erfassung
der Arbeitszeit.

\subsection{Kommunikations- und Toolplanung}
Ein weiterer Aspekt ist die Kommunikation in der Gruppe. Wie wurde diese
organisiert? Gab es regelm"a"sige Treffen oder nur ein Kick-Off Meeting
und dann nur eMail basierte Kommunikation? Erfolgte die Kommunikation nur
"uber eMail, oder wurden Skype, ICQ, Foren oder anderen Kommunikationsmittel
eingesetzt? Wurde ein interner Gruppenleiter bestimmt und der Kontakt zur
Gruppe erfolgte "uber ihn, oder gab es direkte Kommunikation mit dem
Projektmanager? Wie setzten sich die Mitglieder der Gruppe untereinander
in Verbindung? Diese und weitere Fragen zur Kommunikation in der Gruppe
sollten in diesem Kapitel behandelt werden.

Das Ziel des Kommunikationsmanagement ist es im Rahmen des Projekts
rechtzeitig und angemessen die Erzeugung, Sammlung, Verbreitung, Speicherung
und L"oschung von Projektinformationen zu gew"ahrleisten.

Hauptprozesse \cite{SCHW09} sind:
\begin{itemize}
  \item Projektbeteiligte identifizieren,
  \item Planung der Kommunikation,
  \item Verbreitung von Informationen,
  \item Steuerung der Erwartungen der Projektbeteiligten,
  \item Bericht "uber die Leistungen.
\end{itemize}

\subsection{Urlaub und Krankheit}
Im Anschluss zu dem Kommunikationsaspekt soll auch die Verf"ugbarkeit der
Projektteilnehmer festgehalten werden.  Die Teilnehmer k"onnten Urlaub haben
oder "uber einen l"angeren Zeitraum krank werden und somit f"ur die Arbeit
nicht zur Verf"ugung stehen. Diese Details sollten ebenso in diesem Abschnitt
dokumentiert werden.

\section{Anhang}
In dem Anhang befinden sich alle angefertigten Unterlagen und Diagramme
welche vom Projektmanager w"ahrend der Projektlaufzeit zusammengestellt
wurden. Des Weiteren k"onnen die Ausgaben der Projektmanagement Software
hier mit aufgef"uhrt werden.

Unter Anderem geh"oren auch die Protokolle, Stundenzettel, Gantt Chart
Diagramme, Arbeitspaket Beschreibungen der jeweiligen Woche und weitere
Unterlagen auch in diesen Abschnitt.

%%%%%%%%%%%%%%%%%%%%%%%%%%%%%%%%%%%%%%%%%%%%%%%%%%%%%%%%%%%%%%%%%%%%%%%%%%%%%%%%
% dokument: bibliographie
%%%%%%%%%%%%%%%%%%%%%%%%%%%%%%%%%%%%%%%%%%%%%%%%%%%%%%%%%%%%%%%%%%%%%%%%%%%%%%%%

\clearpage
\phantomsection
\addcontentsline{toc}{section}{Literatur}
\bibliography{swe1-projekt}



\end{document}